\documentstyle[12pt,fullpage]{article}

\newcommand{\tweak}{{\rm Tweak}}
\newcommand{\sipe}{{\rm Sipe}}
\newcommand{\noah}{{\rm Noah}}
\newcommand{\nonlin}{{\rm Nonlin}}
\newcommand{\abtweak}{{\rm AbTweak}}

\title{\abtweak User's Manual}

\author{Qiang Yang \and Steve Woods}

\date{}

\begin{document}

\maketitle
\begin{center}
Computer Science Department \\ University of Waterloo \\ Waterloo,
Ontario, Canada, N2L 3G1\\
\end{center}

\section*{Acknowledgement}
Support for the research is made available by NSERC operating grant
\#OGP0089686 to Qiang Yang.

\section*{Inquiries}
Please send to :

Steven Woods woods@drev.dnd.ca\\
Defense Research Establishment Valcartier\\
Courcelette, Quebec, Canada\\
G0A 1R0.

or

Qiang Yang qyang@logos.waterloo.edu\\
Department of Computer Science\\
University of Waterloo\\
Waterloo, Ontario, Canada\\
N2L 3G1.

\newpage

\section{Introduction}
\abtweak\ is a nonlinear, hierarchical planner.  
The planning language is the first order predicate logic without
quantifiers and logical connectives.  Given an initial state
description and a goal state description, represented as sets of
literals, it finds a partially ordered, least-commtted set of
operators to achieve the goal.  It is correct in that every solution
plan found necessarily solves the goal from the initial state, and
complete in that it will terminate with a solution if one exists.

\abtweak\ is built on top of the planner Tweak (designed by Chapman).  For
reference of both planners, please see the papers by Chapman
\cite{tweak}, by Qiang Yang and Josh Tenenberg \cite{abtweak}, and the
Master Thesis by Steve Woods \cite{thesis}.

The planner is designed by Qiang Yang and Josh Tenenberg, and
implemented by Steve Woods and Qiang Yang at University of Waterloo.
It runs on either Kyoto or Allegro Common Lisp.

An additional function has been added to Abtweak for testing and
comparison. In Proceedings of AAAI91, McAllester and Rosenblitt
proposed a version of nonlinear planning, called {\em Systematic
Nonlinear Planning}.  We have implemented a version of their planner, 
as well as its abstract version.  The central theme of the new planner
is to use establishment relation (or causal link) to protect
everything achieved so far.  

Below, we give a brief introduction to the basic components of the
planner.  Then we will describe how the planner is used in detail.

\section{Components of \tweak}
\tweak is implemented in four main independent sections: plan
represetation, search, constraint management/plan maintenance, and
modal truth criterion (MTC).

\subsection{Plan and Operator Representation}
Each plan is a {\em lisp structure} with the following slots:
\begin{enumerate}
\item plan-a:  a set of operator structures.
\item 
plan-b: a set of pairs of operators, each pair represents a temporal
order.
\item plan-nc: a set of pairs of variables and constants, each pair being a non-codesignation constraint.
\item plan-cr: a set of establishment relations to be protected.
\end{enumerate}

It is assumed that a variable is anything of the form \$*, where * can
be any character or string.  Any symbols not preceded by "\$" is a
constant.  Unique name assumption applies to constants.

Each operator template is also a lisp structure.  It has the following
slots:
\begin{enumerate}
\item operator-opid: identifier for the operator;
\item operator-name: a s-expression as the name of the operator;
\item operator-cost: a cost value for the operator.  This is static;
\item operator-preconditions: a list of precondition literals for the
operator;
\item operator-effects: a list of effect literals for the operator;
\item (Optional): operator-primary-effects: a list of primary effects
of the operator, used for selecting the operator to achieve a goal.
\end{enumerate}

\subsection{Search}
The search routine is an implementation of the common A* graphsearch
found in Nilsson's book \cite{nilsson}.  Each node in the search space
is a list, which consists of the following information:
\begin{enumerate}
\item the total cost (cost-so-far + heuristic) of the node;
\item the plan-structure;
\item the cost of the plan by summing up the costs of the operators;
\item the parent node of the current node.
\end{enumerate}

The open list is a list of sublists, each containing a list of plans
with the same costs.  Functions for accessing this list is provided,
which will be discussed below.

No attempt is made to check cycles in the search space, as this is
still a difficult open problem.  The user can specify various kinds of
heuristics for selecting the next node in the open list.  This is
described in detail later.

\subsection{Constraint Management}
The constraint management portion of Tweak handles the insertion of
operators, ordering constraints between pairs of operators as well as
the codesignation and non-codesignation of variables and constants
within the operators present in a plan.

Codesignations of two variables and of a variable and a constant 
are handled by simply replacing all variable occurences
in the plan with a single value as codesignation occurs.  For example,
if a plan contained the variable \$var, and a constraint was added
such that \$var was to codesignate with the constant 'block', then all
occurences of
\$var in any plan operator, or other plan structure would be replaced by
'block'.  Similarly, if two variables are made to codesignate, all
occurences of the second are replaced by the first.  Noncodesignation
of \$x and \$y is represented as a pair $(\$x, \$y)$.  The plan has a
slot, plan-nc, which is a list of such pairs,

\subsection{Modal Truth Criterion}
The modal truth criterion (mtc) is a function, designed by Chapman, for
checking whether or not a plan  is correct.  To do this, it takes
$O(n^4)$, where $n$ is the total number of operators in the plan.

In the search routine, (mtc plan) is a function for determining whether the
current plan is a necessarily correct solution.

Note:  white knights are {\em not} used in the \tweak\ implementation,
as we don't believe that it offers computational advantage.  It is
trivial to prove the completeness of the planner without white knights.

\section{How to Use the Planner}

\subsection{Defining your Domain}

The first step in using the planner is to define a domain definition.
For example, if we consider the 2 ring tower of hanoi domain, there
can be two operators, defined as follows:

\begin{verbatim}
(setq op1 (create-operator-instance 
	:opid 'moveb 	 
	:name '(moveb $x $y) 	 
	:preconditions '((not ons $x) (onb $x) (not ons $y) 
		(ispeg $x) (ispeg $y) ) 	 
	:effects '((not onb $x) (onb $y))))

(setq op2 (create-operator-instance 	 
	:opid 'moves 	 
	:name '(moves $x $y) 	 
	:preconditions '((ons $x) (ispeg $x) (ispeg $y) )
	:effects '((not ons $x) (ons $y))))

(setq *operators* (list op1 op2))
\end{verbatim}

Each operator has a name, including a list of parameters that the
operator accepts, a list of preconditions specifying exactly the state
the world must be in so as to allow this operator to be applied, and a
list of effects specifying exactly what the results of applying this
operator are.  Negations are allowed, and are specified simply as (not
predicate) where the positive would be (predicate). These operators
are stored in a global variable *operators* so the planner can find
and use them.


An initial state specification might look like:

\begin{tabbing}
(setq \=initial \\
\>	'((ispeg peg1) (ispeg peg2) (ispeg peg3) (onb peg1) (ons\\
\>	peg1) (not onb peg2) (not onb peg3) (not ons peg2) (not ons peg3)))\\
\end{tabbing}

A goal specification for hanoi might look like:
\begin{quote}
(setq goal '((onb peg3) (ons peg3)))
\end{quote}

\subsection{Directory Structures}
The file, plan.lsp, contains the source of the main routine.  Through
the key words set in this routine, one can invoke either \abtweak, or
\tweak.  \abtweak\ related routines are in the directory Ab-routines/.
\tweak related routines are in Tw-routines/.  A* search related
routines are stored in Search-routines/.  In addition, the sub-directory
Plan-routines/ contains general routines connecting the planners.

\subsection{How to Plan}
The main command for planning appears below:
\begin{tabbing}
(defun plan (initial goal \&key \= \\
			     
\>                                 (planner-mode 'tweak)\\
\>				   (continue-p nil)\\
\>                                 (debug-mode      nil)\\
\>                                 (heuristic-mode 'num-of-unsat-goals)\\
\>                                 (expand-bound      500)\\
\>                                 (generate-bound   1000)\\
\>                                 (open-bound       1000)\\
\>  				   (output-file nil)\\
\>		 		   (cpu-sec-limit 60)   ; 1 min\\
~~;AbTweak related keywords\\
\>                                 (use-primary-effect-p  nil)\\
\>                                 (mp-mode  t)\\
\>                                 (abstract-goal-mode  t) \\
\>                                 (drp-mode         nil)\\
\>                                 (left-wedge-mode t)\\
\>                                 (subgoal-determine-mode  'stack)\\
\>				   (solution-limit 100)\\
\>               )
\end{tabbing}

To run the planner, be sure one is located in the directory Abtweak/.
Invoke lisp interpreter, and load ``init.lsp''.

If $initial$ contains the initial state literals, and $goal$ the goal
state literals, then the following call invokes the planner tweak.
\begin{tabbing}
~~(plan initial goal :planner-mode 'tweak)
\end{tabbing}

Upon termination, the total number of nodes expanded and generated are
returned. Also, the solution plan is stored in a global variable
*solution*.

Some important key words that are related to \tweak\ are briefly
described below:
\begin{description}
\item[planner-mode]  Can be 'tweak, 'abtweak, 'mr, or 'mr-crit.

If 'abtweak is used, then be sure that a list of criticality
assignments are given in *critical-list*.  Also, if left-wedge-mode is
chosen, then a list of level-preference should be given through *k-list*.

In addition, two other planners can be chosen with either $'mr$ or
$\mbox{'mr-crit}$ as the keyword value.  Option 'mr invokes McAllester and
Rosenblitt's systematic nonlinear planner, and mr-crit is an abstract
version of mr.

\item[continue-p] once set, the next lowest cost solution is planned
for, taking the existing open list as the current state space.

\item[heuristic-mode] With 'num-of-unsat-goals, the total number of
unsatisfied goals are taken as a heuristic.  However, if set to be
'user-defined, then the user defined heuristic function (explained
below) is used to evaluate the heuristic.

\item[expand-bound] If the total number of nodes expanded exceeds this
number, then the planner terminates.

\item[generate-bound] If the total number of nodes generated exceeds this
number, then the planner terminates.

\item[open-bound] If the total number of nodes in the open list exceeds this
number, then the planner terminates.

\item[cpu-sec-limit] If the total number of cpu time in seconds
exceeds this number, then planner terminates.

\item[debug-mode] If t, the planner breaks every time a new node is
generated.  Both the parent and the new node are displayed for
debugging purposes.

\item[output-file]  If non-nil, the output is written into the
filename specified.  For example, :output-file ``file1.data''
specifies that the output of the planner should be written into
file1.data. 

\item[solution-limit]  Maximum depth of the search tree, in terms of
the total number of operators in a solution.

\item[use-primary-effect-p] T iff one wants to use primary effects for
goal achievement.  Before using this option, make sure that each
operator in a domain has a set of primary effects specified, besides
the set of effects.  The primary effects are used for backwards
chaining during goal achievement, but the effects are generally used
for checking conflicts and interactions.  An example where primary
effects are used is given in Domains/simple-robot-1.lsp.


\item[Other flags will be described in detail, in later versions of
the manual.]
\end{description}



\subsubsection{Specifying Search Heuristic}
The user can specify a user-defined heuristic function.  This function
value will be added to the cost-so-far, and used to select a node from
the open list.  To activate this, the key-word ``heuristic-mode'' has
to be set true.

The function (user-heuristic) returns a function, which takes a plan
as a parameter, and which evaluates to be a number.  An example of one
such function is as the following:
\begin{tabbing}
(defun {\rm user-heuristic} () \\
~~ `(lambda (plan) ({\rm num-of-unsat-goals} plan)))\\

(defun {\rm num-of-unsat-goals} (plan)\\
~~(declare \\
~~~~(type array plan) )\\
~~  (count-if-not\\
~~  \#'(lambda (ith-goal)\\
~~~~       (hold-p plan 'g ith-goal))\\
~~   (get-preconditions-of-opid 'g plan)))
\end{tabbing}
This function evaluates the total number of unsatisfied goals in the
plan, and use that number as the heuristic value.

\section{AbTweak}
AbTweak combines the advantages of Tweak and Abstrips
\cite{abstrips}, by planning at multiple levels of abstraction.  The
abstraction hierarchy is specified via a criticality assignment to the
literals in the domain.  At each level $i$ of abstraction, all
preconditions of an operator which are lower than $i$ are removed,
resulting in a simpler operator set on that level.  Planning is first
done at the highest level of abstraction, using the operator set on
the corresponding level.  When a correct solution is found at level
$i$ of abstraction, it is {\em refined} to the $i-1$ level by adding
back all preconditions of criticality $i-1$.  It terminates when MTC
returns True at the lowest level of abstraction.

For more information, please see \cite{abtweak}.

\subsection{Criticality Assignments}
The criticality assignment to the literals is done through a global
variable *critical-list*.  Each element of the list is a sublist of
the form:
\begin{verbatim}
  (crit-value  (p $ $) (not q $))
\end{verbatim}
That is, any literal with a predicate $p$ or $\neg q$, with whatever
argument(s), will have a criticality value = {\em crit-value}.

For the Tower of Hanoi domain, this is specified as follows:
\begin{verbatim}

(setq *critical-list-1* '(
   (3 (ispeg $))
   (2  (not onb $) (onb $) )
   (1  (not onm $) (onm $) )
   (0  (not ons $) (ons $) )
)
\end{verbatim}
That is, $(ispeg \$)$ literals has the highest criticality in this
assignment, while the $(ons \$)$ ones the lowest.
The global variable *critical-list* is set in a domain file.

\subsection{Left-Wedge Control Strategy}
Our experience has been that, for most ``good'' abstraction
hierarchies and most problems, the first abstract solution can be
refined successfully to a lowest level solution.  Although not always
true, this observation can be taken as a heuristic that leads to a
heavier invest in computation time to the first abstract solution than
the rest.  The result, is a complete search control method, called the
{\em left-wedge}.

To make it possible to execute left-wedge, a list of cost-adjustment
values have to be provided.  Let the $i^{th}$ value of the list be $v_i$,
then the cost of a plan at $k-i$th level abstraction is adjusted by $cost(plan)-v_i$.
This list of values is specified through a global variable *left-wedge-list*.

As an example, the left wedge list for Tower of Hanoi domain can be
specified as follows:
\begin{verbatim}
(setq *left-wedge-list* '(0 1 3 7))
\end{verbatim}

\subsection{The Downward Refinement Property}
Some abstraction hierarchy has the property that, every abstraction
solution has a refinement down to the lowest level of abstraction.  In
that case, one does not have to keep all abstraction solutions in the
OPEN list.  As soon as one abstract solution is found to be correct,
the rest can be saved on a stack, *drp-stack*.  This can often lead to
an exponential amount of saving in computation \cite{downward}.

If :drp-mode is on, then in the abstract direction (vertical), a depth
first search is done.  Backtracking to the previous level occur only
when the current level search fails. 

Flag :drp-debug-mode will enter a break loop whenver an abstract
refinement of a solution is found at an abstract level.  The user can
then inspect an abstract solution, before it is passed down to the
next level for refinement.  In Allegro Common Lisp, type :cont to
resume computation.

\subsection{The Monotonic Property}
The monotonic property is a formalization of the idea of {\em
protection intervals} in classical hierarchical planning.  Once a
solution is found at level $i$, all causal relations for the
preconditions above $i-1$ is computated, and protected during planning
at subsequent levels.  This heuristic has been shown to be complete
\cite{abtweak}.

\subsection{The Abstract Planner}
Now we describe how to invoke the abstract planner.

\begin{enumerate}
\item To call AbTweak, type 
\begin{quote}
(plan initial goal :planner-mode 'abtweak)
\end{quote}
\item To use left-wedge strategy, type
\begin{quote}
(plan initial goal :planner-mode 'abtweak :left-wedge-mode t)
\end{quote}
\item To use the monotonic property, type
\begin{quote}
(plan initial goal :planner-mode 'abtweak :mp-mode t)
\end{quote}
\item To use the downward refinement property, type
\begin{quote}
(plan initial goal :planner-mode 'abtweak :drp-mode t)
\end{quote}
\end{enumerate}

\section{Planner Mode}
A global variable *planner-mode* determines the planner routines that
are run,  Tweak, Abtweak, MR, or MR-crit.

\section{Sample Domains}
Several domains are listed in /Abtweak/Domains/.
They include a simple blocksworld domain, a Nilsson's blocksworld
domain, a Towers of Hanoi domain, a register's content exchange 
domain, a simple robot planning domain, a simple transportation
domain, a simple computer hardware domain, a simple biology
domain,  a simple database query optimization domain, and a simple
natural language style generation domain.

The two robot planning domains are adapted from Sacerdoti's Abstrips
domain, where a robot can move a number of boxes around several rooms,
and the rooms have doors that can be opened or closed.  For a picture
of the domain, consult Sacerdoti's paper on Abstrips, AIJ.  To reduce
search, we have implemented a user-defined heuristic that checks
cycles in this domain.  The file that contains this heuristic is in
/Domains/robot-heuristic.lsp.  This heuristic can be used if (a) the
file is loaded into lisp interpreter, and (b) the option
\begin{quote}
(plan initial goal :heuristic-mode 'user-defined ...) 
\end{quote}
is used.

\section{Other Useful Routines}
\begin{description}
\item[Show Plan]
(show-ops  plan-structure)

input: plan structure

output: a pretty print structure of the plan.

\item[Statistics So Far]
(cur)

Outputs the current statistics including the number of nodes expanded
and generated.

Outputs a list of initial states, goal states, and operator templates.

\item[Length of the Open List]
(length-of-open)

outputs the length of the open list.
\item[First and Last Elements In The Open List]
(first-of-open) or (last-of-open)

outputs the first or the last plan structure on the plan list.
\end{description}

\begin{thebibliography}{99}

\bibitem{tweak}
D. Chapman, "Planning for Conjunctive Goals", {\em Artificial
Intelligence}, (32), pages 333-377, 1987.

\bibitem{abtweak}
Q. Yang and J. Tenenberg, "ABTWEAK: Abstracting a Nonlinear, Least
Commitment Planner", Proceedings of the Eighth National Conference on
Artificial Intelligence, pages 199-204, 1990.

\bibitem{downward}
F. Bacchus and Q. Yang, "The Downward Refinement Property," 
Proceedings of the IJCAI 91.

\bibitem{mr}
McAllester and Rosenblitt, {\em Systematic Nonlinear Planning},
Proceedings of the AAAI91.

\bibitem{nilsson}
N. J. Nilsson, {\em Principles of Artificial Intelligence}, Morgan
Kaufmann Publishers, 1980.

\bibitem{abstrips}
Sacerdoti, Earl 1974.
 {\em Planning in a hierarchy of abstraction spaces.}
 {\em Artificial Intelligence} 5:115--135.

\bibitem{thesis}
Steve Woods, Master Thesis: {\em An Implementation and Evaluation of A
Nonlinear, Hierarchical Planner.} Department of Computer Science,
University of Waterloo, 1991.

\bibitem{tech-rep}
Qiang Yang, Josh D. Tenenberg and Steven Woods, ``Abstraction in
Nonlinear Planning,'' Unpublished journal length paper. Available as University
of Waterloo technical report CS-91-65. 52 pages.  A postscript version can be
obtained via anonymous ftp to "cs-archive.uwaterloo.ca''    

\end{thebibliography}

\end{document}
